\documentclass[11pt,a4paper]{letter}
\usepackage[margin=2.5cm]{geometry}
\usepackage{hyperref}
\usepackage{microtype}
\setlength{\parskip}{0.6\baselineskip}
\setlength{\parindent}{0pt}

\signature{Miguel \'{A}ngel D\'{i}az-Campos\\Instituto Nacional de Pediatr\'{i}a\\M\'{e}xico City, M\'{e}xico\\\href{mailto:mdiazc161@unam.edu}{mdiazc161@unam.edu}}
\address{Miguel \'{A}ngel D\'{i}az-Campos\\Instituto Nacional de Pediatr\'{i}a (INP)\\Av.\ Insurgentes Sur 3700, Letra C\\04530 Coyoac\'{a}n, Ciudad de M\'{e}xico\\M\'{e}xico}
\date{}

\begin{document}

\begin{letter}{Editorial Office\\\textit{Nature Communications}\\Springer Nature}
\opening{Dear Editor,}

We are pleased to submit our manuscript entitled \textbf{``Multi-method consensus analysis of targeted \textit{in situ} spatial transcriptomics: a framework with pilot demonstration on human lung adenocarcinoma''} for consideration as a Research Article in \textit{Nature Communications}.

\textbf{What the paper does.} We introduce a multi-tool consensus framework for analysing targeted \textit{in situ} spatial transcriptomics data and demonstrate it on a publicly available Xenium lung adenocarcinoma dataset (268{,}034 cells, 289-gene panel). For every analytical family --- cell-type annotation, spatially variable gene detection, spatial domain inference, cell--cell communication, and pseudotime --- we run at least two independent computational tools and retain only findings concordant across methods. The framework yields 288 of 289 panel genes as consensus spatially variable genes (three-method test); a per-cell consensus subset between Banksy and Novae spatial-domain partitionings (ARI~$=0.15$ at the partition level, 35.4\% high-confidence cells at the per-cell level); 20 of 30 cell--cell communication interactions retained under joint Benjamini--Hochberg FDR (LIANA+) and Bonferroni correction (Spacia); and a two-method pseudotime consensus (Slingshot + diffusion pseudotime, $r = 0.478$) restricted to the macrophage compartment.

\textbf{Why it matters.} The dominant Bonferroni-validated cell--cell communication hub features ADAM17 (sheddase) and MUC1 (oncomucin), spatially co-localised across six independent myeloid and stromal sender cell types converging on tumor cells, with Spacia adjusted $p$-values reaching $1.10 \times 10^{-71}$. Activity across both M1-like and M2-like macrophage states, combined with a non-directional macrophage activation continuum (both M1 and M2 markers increase with pseudotime), supports a model of shedding-mediated, non-checkpoint immunosuppression at the tumor--parenchyma interface --- distinct from the canonical checkpoint-exhaustion paradigm and therapeutically actionable through ADAM17 inhibition.

\textbf{Why \textit{Nature Communications}.} The paper combines a methodological framework with a substantive biological case study. The framework is generic across in situ spatial platforms and gene-panel sizes, and we explicitly document the methodological limits (which tools fail at 289 genes, why, and at what panel size each is expected to operate; Supplementary Note S1). The accompanying open-source pipeline is engineered for portability to the multi-sample, multi-organ cohort that we and others will analyse in the next round of Xenium studies. The combined methodological reach and biological depth fit the cross-disciplinary readership of \textit{Nature Communications} better than a methods-only venue.

\textbf{Transparency.} A reviewer-style audit identified five issues in an earlier framing of this work (Supplementary Note S1.3): a downstream test applied to embeddings cast as integer counts, an implausible trajectory configuration, a magnitude-rank threshold mislabelled as FDR, per-test $p$-values not corrected across cross-tests, and software-version drift between code and manuscript. All five were resolved before submission, and the audit memo is included in full as a supplementary note --- both as a reviewer aid and as a transparency statement that the field, in our view, currently lacks.

\textbf{Honest limitations.} The pilot is a single section from a single patient. Cell-level statistical significance is high but does not substitute for biological replication. We frame this paper as a methods + case-study, with hypothesis-generating biological findings. A multi-sample cohort study using the same framework on a 439-gene panel (telomere biology disorder-associated lung and liver fibrosis) is currently being collected and will be reported separately.

\textbf{Data and code availability.} The dataset is publicly available at 10x Genomics (CG000709 Rev~C). The full pipeline (Snakemake DAG, conda environments, Docker image for Spacia) is available at \url{https://github.com/MiguelDiaz02/xenium-spatial-transcriptomics} under the MIT licence; the version of record will be tagged at acceptance.

The manuscript has not been submitted to or considered by any other journal. All authors have approved the submission and declare no competing interests.

We thank you for your consideration and look forward to your decision.

\closing{Sincerely,}

\end{letter}

\end{document}
